%% ============================================================
%%  nota-bellman.tex — De onde vem a equação de Bellman
%%  Aprendizagem por Reforço · UFU · Prof. Marcelo Keese Albertini
%%  Compilar com XeLaTeX (2x)
%% ============================================================
\documentclass[10pt, a4paper]{article}

\usepackage{fontspec}
\usepackage[margin=2.1cm, top=1.5cm, bottom=1.8cm]{geometry}
\usepackage{babel}
\babelprovide[import, main]{brazilian}
\usepackage{mathtools}
\usepackage{amssymb}
\usepackage{amsthm}
\usepackage[most]{tcolorbox}
\usepackage{tikz}
\usepackage{booktabs}
\usepackage{enumitem}
\usepackage{array}
\usepackage{fancyhdr}
\usepackage[hidelinks]{hyperref}

%% ---------------- cores do curso ------------------------------
\definecolor{rlAzul}{HTML}{16324F}
\definecolor{rlAzulClaro}{HTML}{2E5F8A}
\definecolor{rlTeal}{HTML}{18808D}
\definecolor{rlLaranja}{HTML}{E8590C}
\definecolor{rlAmbar}{HTML}{F2A93B}
\definecolor{rlVermelho}{HTML}{C1121F}
\definecolor{rlVerde}{HTML}{2B9348}
\definecolor{rlCinza}{HTML}{5C677D}
\definecolor{rlFundo}{HTML}{FAFBFC}

%% ---------------- tipografia ----------------------------------
\IfFontExistsTF{Fira Sans}{%
  \setsansfont{Fira Sans}[
    ItalicFont = Fira Sans Italic,
    BoldFont = Fira Sans SemiBold,
    BoldItalicFont = Fira Sans SemiBold Italic]
}{\setsansfont{TeX Gyre Heros}}
\IfFontExistsTF{Fira Mono}{\setmonofont{Fira Mono}[Scale=0.88]}%
  {\setmonofont{Latin Modern Mono}[Scale=0.9]}
\renewcommand{\familydefault}{\sfdefault}

\setlength{\parindent}{0pt}
\setlength{\parskip}{0.42em}
\linespread{1.04}

%% ---------------- cabeçalho / rodapé --------------------------
\pagestyle{fancy}
\fancyhf{}
\renewcommand{\headrulewidth}{0pt}
\renewcommand{\footrulewidth}{0pt}
\fancyfoot[L]{\footnotesize\color{rlCinza}Aprendizagem por Reforço
  \textbullet\ Nota técnica \textbullet\ Equação de Bellman}
\fancyfoot[R]{\footnotesize\color{rlCinza}\thepage\,/\,\pageref{fim}}

%% ---------------- seções --------------------------------------
\newcounter{secao}
\newcommand{\secao}[1]{%
  \vspace{0.8em}%
  \stepcounter{secao}%
  \noindent{\color{rlLaranja}\rule{1.6em}{2pt}}\hspace{0.5em}%
  {\color{rlCinza}\footnotesize\bfseries\thesecao}\hspace{0.35em}%
  {\color{rlAzul}\large\bfseries #1}%
  \par\vspace{-0.15em}%
  {\color{rlAzul!22}\rule{\linewidth}{0.6pt}}%
  \par\vspace{0.3em}}

%% ---------------- caixas --------------------------------------
\tcbset{ncaixa/.style={
  enhanced, boxrule=0pt, frame hidden, arc=1.2mm,
  left=2.4mm, right=2.4mm, top=1.3mm, bottom=1.5mm, boxsep=0.7mm,
  before skip=0.75em, after skip=0.75em,
  borderline west={2.2pt}{0pt}{#1},
  colback=rlFundo, fonttitle=\bfseries\small, coltitle=#1,
  attach title to upper={\par\vspace{0.3ex}}, breakable,
}}
\newtcolorbox{defbox}[2][]{ncaixa=rlTeal, title={#2}, #1}
\newtcolorbox{ideiabox}[1][]{ncaixa=rlAmbar, title={Leitura},
  colback=rlAmbar!9, coltitle=rlLaranja, #1}
\newtcolorbox{alertabox}[2][]{ncaixa=rlVermelho, title={#2},
  colback=rlVermelho!5, #1}
\newtcolorbox{chave}[1][]{ncaixa=rlVerde, title={Resultado},
  colback=rlVerde!6, #1}

%% caixa de um passo da dedução
\newcounter{passo}
\newtcolorbox{passo}[1]{enhanced, breakable, arc=1.2mm, boxrule=0.7pt,
  colframe=rlAzul!16, colback=white,
  left=2.6mm, right=2.6mm, top=1.6mm, bottom=1.8mm, boxsep=0.6mm,
  before skip=0.8em, after skip=0.8em,
  fonttitle=\bfseries\small, coltitle=rlAzul,
  attach boxed title to top left={xshift=3mm, yshift=-2.4mm},
  boxed title style={colback=rlAzul, arc=0.8mm, boxrule=0pt,
    left=1.6mm, right=1.6mm, top=0.6mm, bottom=0.6mm},
  coltitle=white,
  title={Passo \stepcounter{passo}\thepasso\ \ \textbullet\ \ #1}}

%% ---------------- macros do curso -----------------------------
\newcommand{\E}{\mathbb{E}}
\newcommand{\Prob}{\mathbb{P}}
\newcommand{\R}{\mathbb{R}}
\newcommand{\gS}{\mathcal{S}}
\newcommand{\tran}[2]{P(#1\mid #2)}
\newcommand{\rec}{R}
\DeclarePairedDelimiter{\abs}{\lvert}{\rvert}
\DeclarePairedDelimiter{\norm}{\lVert}{\rVert}
\newcommand{\norminf}[1]{\norm{#1}_{\infty}}
\newcommand{\um}{\mathbf{1}}

\hyphenation{apren-di-za-gem re-com-pen-sa se-quen-cial ho-mo-ge-nei-da-de
  geo-mé-tri-co es-pe-ran-ça}

%% ============================================================
\begin{document}

%% ---------------- faixa de título -----------------------------
\begin{tikzpicture}[remember picture, overlay]
  \fill[rlAzul] ([yshift=0cm]current page.north west)
    rectangle ([yshift=-2.95cm]current page.north east);
  \node[anchor=north west, text width=0.70\paperwidth, inner sep=0pt]
    at ([xshift=2.1cm, yshift=-0.75cm]current page.north west) {%
      {\color{rlAmbar}\footnotesize\bfseries
       UFU \ \textbullet\ \ APRENDIZAGEM POR REFORÇO \ \textbullet\ \ NOTA TÉCNICA}\\[0.4em]
      {\color{white}\Large\bfseries De onde vem a equação de Bellman}\\[0.3em]
      {\color{white!72!rlAzul}\small Dedução a partir da definição de um
       processo de recompensa de Markov}};
  \node[anchor=north east, text width=4.6cm, align=right, inner sep=0pt]
    at ([xshift=-2.1cm, yshift=-0.82cm]current page.north east) {%
      {\color{white!70!rlAzul}\scriptsize Prof.\ Marcelo Keese Albertini\\[0.15em]
       complemento da Aula 1}};
\end{tikzpicture}
\vspace*{1.35cm}

A pergunta que esta nota responde: por que uma esperança sobre \emph{infinitas
trajetórias de comprimento infinito} pode ser reescrita como um sistema de
$\abs{\gS}$ equações lineares? A resposta não é um truque só — são cinco passos,
cada um usando uma ferramenta diferente. Vale separá-los, porque cada hipótese
do modelo entra em um passo específico, e saber \emph{onde} ela entra é saber o
que quebra quando ela cai.

%% ============================================================
\secao{Ponto de partida}

\begin{defbox}{O que estamos supondo}
Um \textbf{processo de recompensa de Markov} $(\gS, P, \rec, \gamma)$ com:
\begin{itemize}[leftmargin=1.2em, itemsep=0.15ex, topsep=0.25ex]
  \item $\gS$ finito, e transições $\tran{s'}{s} = \Prob(s_{t+1}=s'\mid s_t=s)$
        que \textbf{não dependem de $t$} (homogeneidade temporal);
  \item recompensa esperada $\rec(s) = \E[r_t \mid s_t = s]$, também
        independente de $t$, e \textbf{limitada}:
        \mbox{$\abs{\rec(s)} \le \rec_{\max} < \infty$};
  \item horizonte \textbf{infinito} e desconto $\gamma \in [0,1)$.
\end{itemize}
\end{defbox}

As duas definições que vamos manipular, e nada além delas:
\[
  G_t \;=\; \sum_{k=0}^{\infty} \gamma^{k} r_{t+k},
  \qquad\qquad
  V(s) \;=\; \E\!\left[\, G_t \;\middle|\; s_t = s \,\right].
\]

$G_t$ é uma \emph{variável aleatória} — muda a cada trajetória sorteada.
$V(s)$ é um \emph{número} por estado. O objetivo é chegar em
$V(s) = \rec(s) + \gamma \sum_{s'} \tran{s'}{s}\,V(s')$.

%% ============================================================
\secao{A dedução, passo a passo}

\begin{passo}{álgebra pura}
\textbf{O retorno é auto-similar.} Separe a primeira parcela e fatore $\gamma$
do que sobra:
\[
  G_t = r_t + \gamma r_{t+1} + \gamma^2 r_{t+2} + \cdots
      = r_t + \gamma\underbrace{\bigl(r_{t+1} + \gamma r_{t+2} + \cdots\bigr)}
        _{\textstyle = \; G_{t+1}}
  \qquad\Longrightarrow\qquad
  \boxed{\,G_t = r_t + \gamma\,G_{t+1}\,}
\]
Nenhuma probabilidade foi usada aqui: isto vale trajetória a trajetória. É a
observação que faz todo o resto funcionar — a \textbf{cauda} da soma, depois de
fatorar $\gamma$, é uma cópia exata da soma original deslocada um passo.
\end{passo}

\begin{alertabox}{Por que o desconto tem de ser geométrico}
Suponha pesos quaisquer, $G_t = \sum_{k\ge0} w_k\, r_{t+k}$. Para que a cauda
seja proporcional a $G_{t+1}$, precisamos de uma constante $c$ com
$\sum_{k\ge1} w_k r_{t+k} = c\sum_{k\ge0} w_k r_{t+1+k}$ para \emph{toda}
sequência de recompensas — ou seja, $w_{k+1} = c\,w_k$ para todo $k$, o que
força $w_k = w_0\,c^{k}$.

O desconto geométrico não é apenas \emph{suficiente} para haver recursão
estacionária: é a \emph{única} ponderação que a produz. Com $\gamma^k$ o
problema é auto-similar; com qualquer outro perfil de pesos, não existe equação
de Bellman.
\end{alertabox}

\begin{passo}{linearidade da esperança}
\textbf{Passe a esperança para dentro.} Condicionando a $s_t = s$ e usando
$G_t = r_t + \gamma G_{t+1}$:
\[
  V(s) \;=\; \E[\,r_t \mid s_t = s\,]
        \;+\; \gamma\,\E[\,G_{t+1} \mid s_t = s\,]
        \;=\; \rec(s) \;+\; \gamma\,\E[\,G_{t+1} \mid s_t = s\,],
\]
onde o primeiro termo é exatamente a definição de $\rec(s)$ no MRP.

\end{passo}

\begin{passo}{lei da esperança total}
\textbf{Abra o futuro pelo próximo estado.} O termo difícil é
$\E[G_{t+1}\mid s_t=s]$ — uma esperança sobre todo o futuro, condicionada ao
presente. Particione pelos valores possíveis de $s_{t+1}$:
\[
  \E[\,G_{t+1} \mid s_t = s\,]
  \;=\; \sum_{s' \in \gS} \tran{s'}{s}\;
        \E[\,G_{t+1} \mid s_t = s,\; s_{t+1} = s'\,].
\]
Esta identidade é sempre verdadeira, para qualquer processo estocástico. Ela
não usa Markov nem homogeneidade — apenas reorganiza a esperança.
\end{passo}

\begin{passo}{propriedade de Markov}
\textbf{Esqueça de onde você veio.} $G_{t+1}$ é função apenas de
$r_{t+1}, r_{t+2}, \ldots$, que por sua vez dependem de
$s_{t+1}, s_{t+2}, \ldots$ A propriedade de Markov diz que, \emph{dado}
$s_{t+1}$, essa trajetória futura é independente de $s_t$ e de toda a história
anterior. Logo o condicionamento em $s_t = s$ é redundante e pode ser
descartado:
\[
  \E[\,G_{t+1} \mid s_t = s,\; s_{t+1} = s'\,]
  \;=\; \E[\,G_{t+1} \mid s_{t+1} = s'\,].
\]
\end{passo}

\begin{passo}{homogeneidade temporal + horizonte infinito}
\textbf{Reconheça o mesmo objeto.} Como $P$ e $\rec$ não dependem de $t$, e
como de $t+1$ em diante ainda resta um futuro infinito — exatamente como
restava em $t$ — a esperança acima é a \emph{mesma função} $V$ aplicada em
$s'$:
\[
  \E[\,G_{t+1} \mid s_{t+1} = s'\,] \;=\; V(s').
\]
\end{passo}

Encadeando os cinco passos:

\begin{chave}
\[
  V(s) \;=\; \underbrace{\rec(s)}_{\text{recompensa imediata}}
       \;+\; \gamma \sum_{s' \in \gS} \tran{s'}{s}\; V(s')
  \qquad\text{para todo } s \in \gS.
\]
Em forma vetorial, com $V, \rec \in \R^{N}$ e $P$ a matriz $N\times N$ de
transições: $\;V = \rec + \gamma P V$.
\end{chave}

%% ============================================================
\secao{Onde cada hipótese entrou}

\begin{center}
\footnotesize
\setlength{\tabcolsep}{6pt}
\renewcommand{\arraystretch}{1.22}
\begin{tabular}{@{}c l l@{}}
  \toprule
  \textbf{Passo} & \textbf{Ferramenta} & \textbf{Hipótese consumida} \\
  \midrule
  1 & álgebra elementar        & desconto \emph{geométrico} \\
  2 & linearidade da esperança & $\gamma<1$ e recompensa limitada \\
  3 & lei da esperança total   & \textcolor{rlCinza}{nenhuma — sempre válida} \\
  4 & ---                      & propriedade de \emph{Markov} \\
  5 & ---                      & homogeneidade temporal + horizonte infinito \\
  \bottomrule
\end{tabular}
\end{center}

\begin{ideiabox}
Repare que Markov aparece \textbf{uma única vez}, e não onde a intuição sugere.
Ele não é o que permite a recursão — quem faz isso é o desconto geométrico, no
Passo 1. Markov é o que permite que a recursão seja escrita em termos de
\emph{estados} em vez de \emph{histórias}. Sem ele a equação ainda existiria,
na forma $V(h_t) = \rec(h_t) + \gamma\sum_{h'} P(h'\mid h_t) V(h')$, mas com um
domínio que cresce exponencialmente com $t$ — inútil na prática.
\end{ideiabox}

%% ============================================================
\secao{O que quebra quando cada hipótese cai}

\textbf{Desconto não geométrico.} Sem $w_k = \gamma^k$, a cauda não é
proporcional a $G_{t+1}$ e o Passo 1 falha. Não há recursão estacionária, logo
não há ponto fixo a resolver.

\textbf{Estado não é de Markov.} O Passo 4 falha e sobra
$\E[G_{t+1}\mid s_t, s_{t+1}]$, que ainda depende de onde o processo veio.
A função valor passa a ser indexada por histórias. É precisamente por isso que
a escolha da representação de estado é uma decisão de modelagem tão crítica.

\textbf{Horizonte finito.} O Passo 5 falha: com $n$ passos restantes o futuro
não é ``igual'' ao presente. Escrevendo $V^{(n)}(s)$ para o valor com $n$ passos
pela frente, o que se obtém não é uma equação de ponto fixo, mas uma
\emph{recursão entre objetos diferentes}:
\[
  V^{(0)}(s) = 0,
  \qquad
  V^{(n)}(s) = \rec(s) + \gamma \sum_{s'} \tran{s'}{s}\, V^{(n-1)}(s').
\]

\begin{ideiabox}
Esta é, letra por letra, a iteração de programação dinâmica vista em aula. A
leitura vale ouro: \textbf{$V_k$ da iteração não é ``uma aproximação de $V$''
no sentido vago — é o valor exato do problema com horizonte $k$}. Convergir
para $V$ significa que o problema de horizonte $k$ se aproxima do de horizonte
infinito, e $\gamma^k$ mede exatamente o quanto ainda falta.
\end{ideiabox}

\textbf{$\gamma = 1$ com horizonte infinito.} Duas falhas simultâneas, que são a
mesma vista de dois ângulos. Do lado probabilístico, o Passo 2 perde a
convergência absoluta: $G_t$ pode nem estar definido. Do lado algébrico,
$P\um = \um$ (as linhas somam 1), isto é, $1$ é autovalor de $P$; então
$(I - P)\um = 0$ e a matriz $(I-\gamma P)$ é singular. A equação deixa de ter
solução única.

%% ============================================================
\secao{Duas objeções que costumam aparecer}

\textbf{``$V$ aparece dos dois lados — isso não é circular?''}

Não, porque a equação de Bellman \emph{não é a definição} de $V$. A definição é
$V(s) = \E[G_t \mid s_t = s]$, que não menciona $V$ à direita. Bellman é uma
\emph{propriedade} que essa função demonstravelmente satisfaz — foi o que os
cinco passos estabeleceram. A circularidade aparente é uma restrição sobre um
objeto já definido, não uma tentativa de defini-lo.

Resta saber se a restrição determina $V$ sozinha. Determina: para $\gamma<1$,
$(I-\gamma P)$ é invertível, então $V = (I-\gamma P)^{-1}\rec$ é a única solução.

\smallskip
\textbf{``Como sei que a solução do sistema é a função valor, e não outra coisa?''}

Porque as duas coincidem explicitamente. Como $\norminf{\gamma P} = \gamma < 1$,
vale a série de Neumann:
\[
  (I - \gamma P)^{-1} = \sum_{k=0}^{\infty} \gamma^{k} P^{k}
  \qquad\Longrightarrow\qquad
  V = \sum_{k=0}^{\infty} \gamma^{k} P^{k} \rec .
\]
E $(P^{k}\rec)(s) = \E[\,r_{t+k} \mid s_t = s\,]$, pois $P^k$ é a matriz de
transição em $k$ passos. Portanto
\[
  V(s) = \sum_{k=0}^{\infty} \gamma^{k}\,\E[\,r_{t+k} \mid s_t = s\,]
       = \E\!\left[\sum_{k=0}^{\infty} \gamma^{k} r_{t+k} \;\middle|\; s_t=s\right]
       = \E[\,G_t \mid s_t = s\,],
\]
que é a definição de onde partimos. O círculo se fecha: a solução do sistema
linear \emph{é} a soma descontada esperada, e não apenas algum ponto fixo
conveniente.

\vfill
{\footnotesize\color{rlCinza}
Complemento da Aula 1 do curso de Aprendizagem por Reforço. Material adaptado e
expandido de \textbf{CS234}, Emma Brunskill, Stanford University.
Prof.\ Marcelo Keese Albertini, Faculdade de Computação, Universidade Federal
de Uberlândia.}

\label{fim}
\end{document}
